O uso de dispositivos vestíveis e de balanças inteligentes ampliou o monitoramento da
composição corporal fora de ambientes clínicos. Entre esses dispositivos estão as
balanças de bioimpedância de oito eletrodos, que estimam de forma segmentada
parâmetros como massa muscular, percentual de gordura, gordura visceral, idade
corporal e metabolismo basal.

Os aplicativos que acompanham essas balanças costumam mostrar apenas os números, sem
interpretá-los nem traduzi-los em recomendações para usuários leigos. A inteligência
artificial (IA) pode preencher essa lacuna, transformando os dados em interpretações
personalizadas e em apoio ao acompanhamento do progresso.

Este trabalho propõe um aplicativo com assistente virtual que interpreta
automaticamente os dados da balança CF577BLE (Figura~\ref{fig:balanca}). A leitura é
feita direto do dispositivo pelo navegador, pela API Web Bluetooth, sem depender de
serviços em nuvem. A ideia inicial era usar um modelo de linguagem de grande porte
(\emph{Large Language Model} --- LLM), mas testes posteriores indicaram que um
\emph{Small Language Model} (SLM) especializado e executado localmente atenderia
melhor. O modelo adotado foi o Qwen 2.5 (7B), que interpreta as métricas de
bioimpedância localmente e mantém os dados privados.

O objetivo é transformar os dados brutos em interpretações claras e recomendações
personalizadas, tornando o acompanhamento da saúde mais útil e seguro para o usuário.

\figura{figuras/balanca}{Balança de bioimpedância modelo CF577BLE, com os oito eletrodos (plataforma e manípulo) e a transmissão por Bluetooth}{fig:balanca}{width=.7\linewidth}

O desenvolvimento do \sintonia\ resultou em uma plataforma funcional que lê a
bioimpedância pela Web Bluetooth, calcula os indicadores corporais com equações
validadas e gera interpretações em linguagem natural. O trabalho mostrou que é viável
usar um \emph{Small Language Model} local, como o Qwen 2.5 (7B), para ajudar o usuário
a entender sua composição corporal sem enviar dados de saúde a serviços externos.
Assim, as informações da balança ficam mais fáceis de interpretar, e o projeto serve
de base para estudos futuros em saúde digital e composição corporal.
